% Source: https://artofproblemsolving.com/community/c3h2110478_help_i_have_one_solution
% Stable identity: aops-topic-2110478-q01
Since $AB=AC$, the vertex $A$ lies on the perpendicular bisector of $BC$, so $A$ is the midpoint of one of the two arcs $BC$. The points $B$ and $C$ split the circle into an arc of $102^\circ$ and an arc of $360^\circ-102^\circ=258^\circ$, and $A$ can lie on either one. The inscribed angle $\angle BAC$ is half of the arc $BC$ that does \emph{not} contain $A$.

\begin{asy}
usepackage("amsmath");
size(7cm);
defaultpen(fontsize(10pt));
picture p1, p2;
// case 1: arc BC not containing A is 102
pair A=dir(90), B=dir(219), C=dir(321);
draw(p1,unitcircle); draw(p1,arc((0,0),1,219,321),red+linewidth(1.6));
draw(p1,A--B--C--cycle,linewidth(0.9)); draw(p1,(0,0)--B,gray+dashed); draw(p1,(0,0)--C,gray+dashed);
dot(p1,A); dot(p1,B); dot(p1,C); dot(p1,(0,0));
label(p1,"$A$",A,N); label(p1,"$B$",B,SW); label(p1,"$C$",C,SE); label(p1,"$O$",(0,0),N);
label(p1,"$102^\circ$",dir(270),S,red); label(p1,"$51^\circ$",A+(0,-0.42),fontsize(8pt));
label(p1,"Case 1",(0,-1.45));
// case 2: arc BC containing A is 102
pair A2=dir(90), B2=dir(141), C2=dir(39);
draw(p2,unitcircle); draw(p2,arc((0,0),1,39,141),red+linewidth(1.6));
draw(p2,A2--B2--C2--cycle,linewidth(0.9)); draw(p2,(0,0)--B2,gray+dashed); draw(p2,(0,0)--C2,gray+dashed);
dot(p2,A2); dot(p2,B2); dot(p2,C2); dot(p2,(0,0));
label(p2,"$A$",A2,N); label(p2,"$B$",B2,W); label(p2,"$C$",C2,E); label(p2,"$O$",(0,0),S);
label(p2,"$102^\circ$",dir(115),NW,red); label(p2,"$129^\circ$",A2+(0,-0.22),fontsize(8pt));
label(p2,"Case 2",(0,-1.45));
add(p1); add(shift(2.8,0)*p2);
label("\mbox{}",(min(currentpicture,user=true)+max(currentpicture,user=true))/2,black); // keeps the text after the figure black when LaTeX typesets the labels inline
\end{asy}

\textbf{Case 1: $A$ is not on the $102^\circ$ arc.} Then $\angle A=\dfrac{102^\circ}{2}=51^\circ$ and $\angle B=\angle C=\dfrac{180^\circ-51^\circ}{2}=64.5^\circ$.

\textbf{Case 2: $A$ is on the $102^\circ$ arc.} Then $\angle A$ subtends the $258^\circ$ arc, so $\angle A=129^\circ$ and $\angle B=\angle C=\dfrac{180^\circ-129^\circ}{2}=25.5^\circ$.

So the angles are either $\boxed{51^\circ,\ 64.5^\circ,\ 64.5^\circ}$ or $\boxed{129^\circ,\ 25.5^\circ,\ 25.5^\circ}$.
